\section{City With the Smallest Number of Neighbors at a Threshold Distance}
\label{sec:graph-city-smallest-neighbors}

\problemstatement{Given $n$ cities connected by weighted edges and a
$\textit{distanceThreshold}$, find the city from which the
\textbf{fewest} other cities are reachable within
$\textit{distanceThreshold}$. If multiple cities tie for fewest,
return the one with the \textbf{largest} index.}

\begin{patternbox}
This closing problem of the chapter is a direct application of
Section~\ref{sec:graph-floyd-warshall}: once \textbf{all-pairs}
shortest distances are known, answering ``how many cities are within
threshold $T$ of city $i$'' for every $i$ is a simple linear scan over
each row of the distance matrix -- no further graph algorithm is
needed.
\end{patternbox}

\subsection*{Reduction}

\begin{enumerate}
  \item Build the $n \times n$ distance matrix and run Floyd-Warshall
        to fill in all-pairs shortest distances.
  \item For each city $i$, count how many cities $j \neq i$ satisfy
        $\textit{dist}[i][j] \le \textit{distanceThreshold}$.
  \item Track the city with the smallest count, breaking ties by
        preferring the \textbf{larger} index (scan $i$ from $0$ to
        $n{-}1$, and use $\le$ rather than $<$ when comparing counts,
        so a later equal-or-better city always overwrites an earlier
        one).
\end{enumerate}

\subsection*{C++ implementation}

\begin{lstlisting}
#include <bits/stdc++.h>
using namespace std;
const int INF = 1e9;

int findTheCity(int n, vector<vector<int>>& edges, int distanceThreshold) {
    vector<vector<int>> dist(n, vector<int>(n, INF));
    for (int i = 0; i < n; i++) dist[i][i] = 0;
    for (auto& e : edges) {
        dist[e[0]][e[1]] = e[2];
        dist[e[1]][e[0]] = e[2];
    }

    for (int k = 0; k < n; k++) {
        for (int i = 0; i < n; i++) {
            for (int j = 0; j < n; j++) {
                if (dist[i][k] < INF && dist[k][j] < INF) {
                    dist[i][j] = min(dist[i][j], dist[i][k] + dist[k][j]);
                }
            }
        }
    }

    int bestCity = -1, minCount = n;
    for (int i = 0; i < n; i++) {
        int count = 0;
        for (int j = 0; j < n; j++) {
            if (i != j && dist[i][j] <= distanceThreshold) count++;
        }
        if (count <= minCount) { // <= prefers the larger index on ties
            minCount = count;
            bestCity = i;
        }
    }
    return bestCity;
}

int main() {
    vector<vector<int>> edges = {{0,1,3},{1,2,1},{1,3,4},{2,3,1}};
    cout << findTheCity(4, edges, 4) << endl; // 3
    return 0;
}
\end{lstlisting}

\subsection*{Dry run}

\dryrun{Take $n=4$, edges $0{-}1(3),1{-}2(1),1{-}3(4),2{-}3(1)$,
$\textit{distanceThreshold}=4$.}

All-pairs distances after Floyd-Warshall: from $0$: $[0,3,4,5]$ (to
$2$: $3+1=4$ via $1$; to $3$: $3+1+1=5$ via $1,2$, shorter than the
direct $1{-}3$ edge route $3+4=7$). From $1$: $[3,0,1,2]$. From $2$:
$[4,1,0,1]$. From $3$: $[5,2,1,0]$ (to $1$: $1+1=2$ via $2$, shorter
than the direct edge $4$). Counts within threshold $4$ (excluding
self): city $0$ sees $1$ (dist $3$) and $2$ (dist $4$) -- $2$ cities
($3$ at distance $5$ is excluded). City $1$ sees all three others --
$3$ cities. City $2$ sees all three others -- $3$ cities. City $3$
sees $1$ (dist $2$) and $2$ (dist $1$) -- $2$ cities ($0$ at distance
$5$ is excluded). Cities $0$ and $3$ tie at $2$ neighbors each; the
tie-break (larger index) selects city $3$.

\begin{complexitybox}
\textbf{Time Complexity.} $O(V^3)$ for Floyd-Warshall, plus $O(V^2)$
for the counting scan: $O(V^3)$ overall.
\textbf{Space Complexity.} $O(V^2)$ for the distance matrix.
\end{complexitybox}

\begin{takeawaybox}
This chapter closes by demonstrating exactly why Floyd-Warshall earns
its $O(V^3)$ cost despite being asymptotically worse than running
Dijkstra from every node ($O(V(V+E)\log V)$ for sparse graphs): once
\emph{every} pairwise distance is available at once, follow-up
questions like this one (counting neighbors within a threshold, for
every city) become trivial linear scans -- exactly the same trade-off
this book has made repeatedly between an upfront, fuller computation
and cheap, flexible queries afterward.
\end{takeawaybox}
